%%%%%%%%%%  scan idx 237  |  folio 224  %%%%%%%%%%
\textbf{Remarque 1.1.2.} Le lemme précédent s'applique notamment au cas où \(i\) est une \uline{immersion régulière} (VII~1.4). En effet, on a localement, grâce au complexe de Koszul, une résolution gauche finie de \(i_*(\underline{O}_X)\) par des \(\underline{O}_Y\)-Modules libres de type fini, donc \(i_*(\underline{O}_X)\) est parfait (I 4.8 et 4.14) (donc a fortiori pseudo-cohérent).

\textbf{Lemme 1.1.3.} \uline{Soit un carré cartésien de schémas}
\[
\begin{tikzcd}
X \arrow[d,"i"'] & X' \arrow[l,"f"'] \arrow[d,"i'"]\\
Y & Y'{,} \arrow[l,"g"']
\end{tikzcd}
\]
\uline{où} \(i\) \uline{est une immersion fermée}, \uline{et soit} \(E\) \uline{un complexe de} \(\underline{O}_X\)-\uline{Modules}. \uline{Alors} la \uline{flèche canonique de changement de base} (EGA~III 1.4.15.4)
\[
  g^*i_*(E)\longrightarrow i'_*f^*(E)
\]
\uline{est un isomorphisme}.

(ii) \uline{Posons} \(E'=f^*(E)\). \uline{Si} \(g\) \uline{est plat et} \(i_*(E)\) \(n\)-\uline{pseudo-cohérent}, \uline{alors} \(i'_*(E')\) \uline{est} \(n\)-\uline{pseudo-cohérent}. \uline{Inversement}, \uline{si} \(g\) \uline{est fidèlement plat et} \(i'_*(E')\) \(n\)-\uline{pseudo-cohérent}, \uline{alors} \(i_*(E)\) \uline{est} \(n\)-\uline{pseudo-cohérent}.

\uline{Preuve}. L'assertion se vérifie trivialement sur les fibres.

